\paragraph{Schwinger--Dyson insertion, measure Jacobians and null directions of the replica functional (conditional; not closure).}
Let \(Z_n(\varepsilon)\) be the replica functional of the paragraph on the replica insertion identity, built from the matrix data of Equation~\eqref{eq:Smatrix} at finite \(N\) and Presentation~A: \(\varepsilon\) deforms a datum of the block or of the induced density and the insertion is \(\mathcal I_n=\partial_\varepsilon S_{E,n}\). Writing the integrand as \(j_\varepsilon\,e^{-S_{E,\varepsilon}}\), with \(j_\varepsilon\) the ratio of the integration measure to the reference matrix measure, differentiation under the integral sign gives, at finite \(N\), the exact identity
\[
  \partial_\varepsilon\ln Z_n=-\langle\mathcal I_n\rangle_n+\langle\partial_\varepsilon\ln j_\varepsilon\rangle_n\,,
\]
and the Schwinger--Dyson (loop) equations of the matrix integral~\cite{EynardOrantin2009}, \(\langle\mathcal V(S_{E,n})\rangle_n=\langle\operatorname{Div}\mathcal V\rangle_n\) for compactly supported vector fields \(\mathcal V\) on the matrix entries, give zero expectation to every equation-of-motion insertion and reduce insertions to contact terms. For the flat Lebesgue matrix measure of Equation~\eqref{eq:Smatrix} the Jacobian term vanishes identically; for Moyal corrections, field redefinitions and the measure of the matrix integral disclaimed in the preceding paragraph it must be retained, and it is not computed anywhere in this book.

Two consequences bound what the loop equations can supply for Problem~B. (i)~Unimodular deformations are null directions: if the \(\varepsilon\)-dependence is generated by a divergence-free field redefinition, then every \(Z_n(\varepsilon)\) is \(\varepsilon\)-independent and \(\langle\mathcal I_n\rangle_n=\langle\partial_\varepsilon\ln j_\varepsilon\rangle_n=0\) for every \(n\). The compensating rotation acts as \(X^a\mapsto e^{tB}X^a\) with \(B^a{}_b=\partial_b\xi^a\) antisymmetric in \(\eta\), hence \(\operatorname{tr}B=0\) and \((\det e^{tB})^{N^2}=1\) over the \(N^2\) matrix entries, and the adjoint flows \(X\mapsto e^{t\Lambda}Xe^{-t\Lambda}\) have \(\operatorname{tr}\operatorname{ad}_\Lambda=0\) (exact checks). The coefficient \(-4\kappa_X\) required in the paragraph on the matrix gauge realisation cannot therefore be produced by the insertion identity along a compensating rotation or an adjoint flow, which sharpens the recorded conclusion that this coefficient is a property of the modular Hamiltonian of the state at \(\Sigma\), or of a local symmetry under which the compensating rotation would be a gauge transformation, and not of the induced density. (ii)~The loop equations do not imply stationarity even at Gaussian order: for the conical-contact family \(A_n=A_1+(\tfrac1n-n)P+\varepsilon B\) of the block quadratic form, with \(P\) supported at the branch locus (the same \(n\)-dependence as the flat-cone coefficient of the paragraph on the replica cone coefficient), the identity above gives exactly
\[
  \partial_\varepsilon S_D=\operatorname{Tr}(G_1BG_1P)-\tfrac12\operatorname{Tr}(G_1B),\qquad G_1=A_1^{-1},
\]
which is nonzero for a generic block deformation \(B\) and vanishes only on the codimension-one hyperplane \(2\operatorname{Tr}(G_1BG_1P)=\operatorname{Tr}(G_1B)\). In this class \(\delta S_D=0\) is a linear condition on the deformation, a linearised block equation together with a contact term at the tip, and not a consequence of any Schwinger--Dyson identity. These statements are for finite \(N\), the Gaussian conical-contact class, the flat matrix measure of Equation~\eqref{eq:Smatrix}, Presentation~A and the fixed-geometry, fixed-measure replica prescription of the paragraph on the replica insertion identity, into which the conical-locus contribution enters through \(S_{E,n}\); they do not give the locality or large-\(N\) factorisation that would turn contact identities into a Ward identity of the block, do not establish fixed-volume stationarity of the tensor block, and do not select the edge term.

\paragraph{Nearest published statements and the reduced-action charge mechanism (conditional; not closure).}
The nearest published analogues compute entanglement of matrix degrees of freedom rather than the response of a block to fixed-volume deformations: target-space entanglement of eigenvalue subregions with qualitative area agreement in the bubbling-geometry matrix model~\cite{TsuchiyaYamashiro2022}; momentum-space entanglement on fuzzy spheres~\cite{KawamotoKuroki2021}; perturbative interaction corrections to entanglement on fuzzy spaces, which are bulk-extensive and do not obey a pure area law~\cite{AlloucheDou2026}; the entanglement of the black 0-brane with its Hawking radiation in the BFSS model~\cite{ChoudhuryLaurenzano2024}; and the \(c=1\) string, where a dilaton-dependent area-like term is expected in the generalised entropy of bulk subregions but no singlet-sector mechanism accounts for it~\cite{Craps2026}. None of these treats fixed-volume variations of a non-conformal block, and none supplies the missing step of Problem~B. On the continuum side, the first law and equilibrium statements for causal diamonds~\cite{JacobsonVisser2019,Svesko2019} and the type-II-algebra realisation of a variant of the entanglement equilibrium hypothesis~\cite{JensenSorceSperanza2023} presuppose an existing metric or a geometric modular flow, an input that a derivation of Assumption~\ref{ass:EE} from Equation~\eqref{eq:Smatrix} must not use. For the edge-term bookkeeping, the published mechanism nearest to the \(r\)-family argument is the charge mechanism of Einstein--aether thermodynamics: the first law of stationary black holes carries an irreducible entropic aether contribution~\cite{ArataLiberatiNeri2026}, and the additional first-law term of misaligned universal horizons is the Noether charge of a genuine local symmetry \(\delta u^a=f\,a^a\) of the reduced, elliptic gauge-fixed action~\cite{ArataIsmailLiberati2026}. The NAMCG situation is the mirror obstruction, not an application of this mechanism: the compensating rotation is an isotropy of the background and not a local rotation of the matrix index, the densities \(L_r\) agree with \(L_0\) only on the stated background class, the term \(Y(\phi,\mathcal L_\xi X)\) does not drop out of the charge, and no local symmetry of the matrix data carries the coefficient \(-4\kappa_X\); the existence of such a symmetry belongs to Problem~D. These identifications record what the published mechanisms do and do not supply for Problem~B; no result is imported into the matrix theory. Problems~B and~D remain open and PROGRAMME-CONDITIONAL.
